\section{Conservation Competency Ontology}

\label{sec:ontology}

\subsection{Design Principles}

The Conservation Competency Ontology (CCO) organizes conservation knowledge into a directed acyclic graph of learning objectives with prerequisite relationships.
We adopt four design principles:

\begin{enumerate}[leftmargin=*]
    \item \textbf{Ecological Validity:} Objectives reflect skills and knowledge demonstrably associated with effective conservation practice, validated through structured interviews with 142 conservation practitioners across 28 countries.
    \item \textbf{Cultural Adaptability:} The ontology separates universal ecological principles from culturally situated knowledge, enabling localization without structural modification.
    \item \textbf{Granularity Balance:} Objectives are sized to be achievable in 10--30 minute learning sessions, supporting mobile and intermittent access patterns common in field settings.
    \item \textbf{Measurability:} Each objective includes assessment criteria specifiable as machine-evaluable rubrics.
\end{enumerate}

\subsection{Domain Structure}

The CCO encompasses six domains, organized hierarchically:

\begin{table}[H]
\centering
\caption{Conservation Competency Ontology domain structure.}
\label{tab:cco_domains}
\begin{tabularx}{\textwidth}{@{}lccX@{}}
\toprule
\textbf{Domain} & \textbf{Objectives} & \textbf{Levels} & \textbf{Description} \\
\midrule
Species Identification & 834 & 5 & Visual, acoustic, and behavioral identification across taxonomic groups \\
Habitat Ecology & 612 & 4 & Ecosystem structure, function, and dynamics \\
Threat Assessment & 487 & 4 & Anthropogenic and natural threats, risk analysis, vulnerability mapping \\
Policy Frameworks & 391 & 3 & International conventions, national legislation, community governance \\
Community Engagement & 298 & 3 & Stakeholder analysis, participatory methods, conflict resolution \\
Field Methodology & 225 & 4 & Survey design, data collection, monitoring protocols \\
\bottomrule
\end{tabularx}
\end{table}

\subsection{Prerequisite Graph}

Prerequisites are modeled as directed edges in the ontology graph.
We define three edge types:

\begin{itemize}[leftmargin=*]
    \item \textbf{Hard prerequisite} ($\rightarrow_H$): Objective $B$ cannot be attempted without demonstrated mastery of $A$.
    \item \textbf{Soft prerequisite} ($\rightarrow_S$): Knowledge of $A$ significantly aids learning of $B$, but $B$ is approachable without $A$.
    \item \textbf{Co-requisite} ($\leftrightarrow_C$): Objectives $A$ and $B$ are mutually reinforcing and optimally learned in parallel.
\end{itemize}

The complete CCO graph contains 2,847 nodes, 4,213 hard prerequisites, 6,891 soft prerequisites, and 1,247 co-requisite pairs.
We verify acyclicity of the hard prerequisite subgraph and compute topological orderings to ensure feasible learning paths exist for all objectives.

\subsection{Formal Representation}

Let $\mathcal{O} = \{o_1, o_2, \ldots, o_N\}$ denote the set of learning objectives, $N = 2847$.
Each objective $o_i$ is characterized by a tuple:

\begin{equation}
    o_i = (d_i, \ell_i, \mathbf{c}_i, \mathcal{A}_i, \tau_i)
\end{equation}

\noindent where $d_i \in \mathcal{D}$ is the domain, $\ell_i \in \{1, \ldots, L_{d_i}\}$ is the proficiency level within domain $d_i$, $\mathbf{c}_i \in \{0,1\}^{|\mathcal{C}|}$ is a binary vector indicating cultural context tags from the set of cultural contexts $\mathcal{C}$, $\mathcal{A}_i$ is a set of assessment specifications, and $\tau_i \in \mathbb{R}^+$ is the estimated median time to mastery in minutes.

% ============================================================
% 4. Adaptive Learning Engine
% ============================================================
